% from ega4-18-fr.tex
\oldpage[IV-4]{192}

\begin{corollary}[19.2.6]
\label{IV.19.2.6}
Supposons, soit que $S$ et $X$ soient localement noethériens, soit que $g$ et $h$ soient localement de présentation finie. Supposons en outre que les fibres $X_s$ et $Y_s$ soient des préschémas réguliers aux points $u(y)$ et $y$ respectivement, et que les morphismes $g$ et $h$ soient plats au voisinage de $u(y)$ et $y$ respectivement. Alors l'immersion $Y\to X$ est transversalement régulière au point $x$.
\end{corollary}

En effet, il résulte de (19.1.1) que l'immersion $Y_s\to X_s$ est régulière au point $x$, et il suffit d'appliquer (19.2.4, $b)$).

\begin{proposition}[19.2.7]
\label{IV.19.2.7}
\begin{enumerate}
  \item[(i)] Toute $S$-immersion ouverte est transversalement régulière relativement à $S$.
  \item[(ii)] Soient $f:Y\to X$ une $S$-immersion, $g:S'\to S$ un morphisme et posons $X'=X_{(S')}$, $Y'=Y_{(S')}$, $f'=f_{(S')}:Y'\to X'$. Si $f$ est une immersion transversalement régulière relativement à $S$, $f'$ est une immersion transversalement régulière relativement à $S'$. La réciproque est vraie si $g$ est fidèlement plat, et si de plus $X$ et $Y$ sont localement de présentation finie sur $S$.
  \item[(iii)] Si $f:Y\to X$ et $g:Z\to Y$ sont des immersions transversalement régulières relativement à $S$, il en est de même de $f\circ g:Z\to X$.
  \item[(iv)] Soient $X$ un $S$-préschéma, $f:Y\to X$, $g:Z\to Y$ deux $S$-immersions. On suppose, soit que $S$ et $X$ soient localement noethériens, soit que $X$, $Y$ et $Z$ soient localement de présentation finie sur $S$. Alors, si $g$ et $f\circ g$ sont transversalement régulières au point $z\in Z$ relativement à $S$, $f$ est transversalement régulière au point $g(z)$ relativement à $S$.
  \item[(v)] Sous les conditions générales de (iv), supposons que $X$ et $Y$ soient $S$-plats, $Y_s$ régulier au point $g(z)$ et $X_s$ régulier au point $g(z)$; alors, si $f\circ g$ est transversalement régulière au point $z$, il en est de même de $g$.
\end{enumerate}
\end{proposition}

L'assertion (i) est immédiate, et la première assertion de (ii) résulte aussitôt de (19.2.3). Pour démontrer la seconde assertion de (ii), appliquons (19.2.4): $X'$ et $Y'$ étant localement de présentation finie sur $S'$, il résulte d'abord de ce critère que $X'$ et $Y'$ sont plats sur $S'$, donc $X$ et $Y$ sont plats sur $S$ (2.5.1); d'autre part, pour tout $s\in S$, si $s'\in S'$ est un point au-dessus de $s$, $Y'_{s'}\to X'_{s'}$ est par hypothèse une immersion régulière, et on déduit donc de (19.1.5, (ii)) qu'il en est de même de $Y_s\to X_s$, puisque ces préschémas sont localement noethériens et qu'il revient alors au même de dire que l'immersion $Y_s\to X_s$ est régulière ou quasi-régulière. Pour prouver (iii), on se ramène au cas où $Z$ est un sous-préschéma fermé de $Y$ et $Y$ un sous-préschéma fermé de $X$, et on forme comme dans (19.1.5, (iii)) un système de générateurs de l'Idéal de $\sh{O}_X$ définissant $Z$. Pour prouver (iv), on remarque que l'hypothèse jointe à (19.2.4) montre d'abord que $X$, $Y$ et $Z$ sont $S$-plats dans un voisinage de $z$, et il suffit en outre, si $s$ est l'image de $z$ dans $S$, de prouver que l'immersion $f_s:Y_s\to X_s$ est régulière au point $z$, sachant qu'il en est de même de $g_s:Z_s\to Y_s$ et de $f_s\circ g_s$; mais cela résulte de (19.1.5, (iv)). Enfin, pour (v), on raisonne de la même manière que pour (iv); l'hypothèse sur $f\circ g$ implique déjà par (19.2.4, $b)$) que $Z$ est $S$-plat dans un voisinage de $z$, et il en est de même de $X$ et $Y$ par hypothèse; on est alors ramené à prouver que $g_s$ est régulière au point $z$, sachant qu'il en est de même de $f_s\circ g_s$, ce qui découle de (19.1.5, (v)), vu les hypothèses faites sur $X_s$ et $Y_s$.

